\scanpage{129}
\mspage{17}
we must have $\gamma_1=0$. Now assume that $i>1$ and that $h_1=\cdots=h_{i-1}=0$. Then $\alpha_{b+i}=\beta_b\gamma_i+\beta_{b+i}$ and by the same argument as above we conclude that $\gamma_i=0$.

Next, it is clear that the product of the elements of $S[x_1]$ of $x_1$-degrees $\lambda_R(t)$ and $\lambda_R(z\mid t)$ which respectively generate $tR$ and $(z\mid t)R$ is an element of $S[x_1]$ of $x_1$-degree $\lambda_R(z)$ and it generates $zR$, and hence without loss of generality we may assume that $t=z$. Then $t=f(x_1)$ where $f(X_1)\in S[X_1]$. Since the line $x_2=\cdots=x\deleted{n}\inserted{{}_n}=0$ is non-tangential to $t$, the coefficient of $x_1^d$ in $f(X_1)$ must be a unit in $S$ and for all $q<d$ the coefficient \inserted{of} $X_1^q$ in $f(X_1)$ must be in $N$. Applying Hensel's Lemma [C\deleted{\unsure{0|O}}\inserted{o}, Theorem 4] to $f(X_1)$ over the complete local domain $(S,N)$ we can write $f(X_1)=g(X_1)h(X_1)$ where $g(X_1)$ is a monic polynomial of degree $d$ in $S[X_1]$ and $h(X_1)$ is a polynomial in $S[X_1]$ whose constant term is a unit in $S$. Then $t=g(x_1)h(x_1)$ and $h(x_1)$ is a unit in $R$. Now by the uniqueness statement proved above we must have $t'=g(x_1)$ and hence $t'$ divides $t$ in $S[x_1]$.

The final statement follows from the facts that the projective plane over any field has at least three points and that in a projective space over an infinite field there exist points lying outside a given hyper-surface.

LEMMA 4. Let $(R,M)$ be a one dimensional local domain with quotient field $K$. Then the integral closure $S$ of $R$ in $K$ is a principal ideal domain with only a finite number of prime ideals $P_1,\ldots,P_n$ $(n>0)$ so
\mspage{128}
